% =====================================================
% 07_discussion.tex — ~0.5 pg
% =====================================================
\todo{honest limitations}

\paragraph{Linear boundaries only.}
The Lie-bracket analysis in \S\ref{sec:theory} treats attribute
directions as constant vector fields on $\W$. For learned non-linear
edits — e.g. EditGAN's part-conditioned edit
vectors~\cite{Ling2021EditGAN} — the full Lie bracket of generally
varying vector fields is non-zero and may dominate the Hessian term.
Extending the curvature prediction to such cases is open.

\paragraph{Forward mode only goes so far.}
For first-order pushforward and HVP we exploit forward mode
maximally. Pure reverse-mode setups (existing Grad-CAM frameworks
trained on classifier loss) cannot directly compete, but
high-dimensional latent-space optimisation (e.g. PTI) is best done in
reverse. We consider forward mode complementary, not a replacement.

\paragraph{Domain coverage and generators.}
We validated on three style-based generators (HiGAN bedroom, StyleGAN2
FFHQ, HiGAN church). Diffusion-based generators have a different
mathematical structure (stochastic, time-conditional) and our
framework as stated does not directly apply; an analogous
score-function-Hessian formulation is open work.

\paragraph{User study scope.}
Our user study covers $\sim 90$ image–attribute pairs across three
domains. Results are statistically significant but cannot speak to
artist-grade editing workflows where context-sensitive judgement
dominates.

\paragraph{Theoretical assumptions.}
We assume $G$ is $C^2$ at the sampled latents. ReLU non-smoothness
forms a measure-zero set in $\W$ and we did not encounter
numerical issues in practice, but a fully rigorous statement requires
piece-wise smooth differential geometry — see supplementary remark.

\paragraph{Beyond GANs.}
The differential-geometric view applies, with adaptation, to any
generative model defined by a smooth map from a latent space. We
expect direct application to VAEs, normalising flows, and
score-based generators after suitable smoothing; this is a clear next
step.
